\documentclass[11pt,a4paper]{article}
\usepackage[margin=1in]{geometry}
\usepackage{amsmath,amssymb}
\usepackage{booktabs}
\usepackage{graphicx}
\usepackage{hyperref}
\usepackage{xcolor}
\usepackage{multirow}
\usepackage{float}

\definecolor{good}{rgb}{0.0,0.5,0.0}
\definecolor{bad}{rgb}{0.7,0.0,0.0}

\title{EEG Scalp-to-In-Ear Signal Prediction:\\
  Iteration 001 --- Baseline Establishment}
\author{Automated Research Loop}
\date{April 7, 2026}

\begin{document}
\maketitle

\begin{abstract}
We present the first iteration of an automated research pipeline for predicting
4-channel in-ear EEG signals from 21-channel scalp EEG (10-20 system). Four
model architectures are compared: closed-form optimal spatial filter, SGD-trained
linear spatial filter, spatio-temporal FIR filter, and nonlinear convolutional
encoder. Using synthetic data (5 subjects, 76,800 samples each, SNR=10\,dB),
the closed-form solution achieves $r=0.887$, SNR$=6.51$\,dB, and coherence$=0.715$,
while all gradient-trained models exhibit negative correlations, indicating
critical training instabilities. Ablation studies confirm that MSE-only loss
outperforms combined spectral/band-power losses. We identify root causes and
propose concrete improvements for Iteration~002.
\end{abstract}

\section{Introduction}

In-ear EEG offers a wearable, unobtrusive alternative to scalp EEG, but
captures a limited spatial view of brain activity. If in-ear signals can be
accurately predicted from scalp recordings, this opens pathways for:
\begin{enumerate}
  \item Transfer learning from large scalp EEG datasets to in-ear models,
  \item Virtual in-ear channel generation for electrode optimization,
  \item Understanding the spatial filtering relationship between montages.
\end{enumerate}

The forward model is: $\mathbf{Y} = \mathbf{M}\mathbf{X} + \boldsymbol{\eta}$,
where $\mathbf{X} \in \mathbb{R}^{21 \times T}$ is scalp EEG,
$\mathbf{M} \in \mathbb{R}^{4 \times 21}$ is the mixing matrix (weighted toward
temporal/parietal channels T7, T8, P7, P8), and $\boldsymbol{\eta}$ is sensor noise.

\section{Methods}

\subsection{Data Generation and Preprocessing}
Synthetic data was generated for 5 subjects with 76,800 samples each at 256\,Hz.
The EEG simulation includes $1/f$ spectral falloff and alpha oscillations.
Preprocessing: bandpass 0.5--45\,Hz, 50\,Hz notch, common average reference,
artifact rejection ($|V| > 150\,\mu$V), z-score normalization, windowing to
256-sample epochs with 50\% overlap. This yielded 2,090 training, 447 validation,
and 449 test windows (chronological split, no temporal leakage).

\subsection{Model Architectures}

\begin{table}[H]
\centering
\caption{Model architectures compared.}
\label{tab:models}
\begin{tabular}{llrl}
\toprule
Model & Type & Parameters & Optimization \\
\midrule
Closed-form & $\mathbf{W}^* = \mathbf{R}_{YX}\mathbf{R}_{XX}^{-1}$ & 84 & Analytic \\
Linear SGD & $\hat{\mathbf{Y}} = \mathbf{W}\mathbf{X}$ & 84 & Adam, 100 epochs \\
FIR Filter & 1D Conv, $L{=}65$ taps & 5,460 & Adam, 200 epochs \\
Conv Encoder & Depthwise-sep + residual & 24,068 & AdamW + Cosine, 300 epochs \\
\bottomrule
\end{tabular}
\end{table}

\subsection{Loss Functions}
The combined loss is $\mathcal{L} = \mathcal{L}_\text{MSE} + \lambda_s \mathcal{L}_\text{spec}
+ \lambda_b \mathcal{L}_\text{band}$, where $\mathcal{L}_\text{spec}$ matches
log-magnitude FFT and $\mathcal{L}_\text{band}$ matches power in 5 EEG bands
($\delta, \theta, \alpha, \beta, \gamma$). Default $\lambda_s = \lambda_b = 0.1$.

\subsection{Evaluation Metrics}
Per-channel Pearson correlation ($r$), RMSE, relative RMSE, SNR (dB),
spectral RMSE, band power correlations (5 bands), and magnitude-squared coherence.

\section{Results}

\subsection{Model Comparison}

\begin{table}[H]
\centering
\caption{Test set performance across all models. Best values in \textbf{bold}.}
\label{tab:results}
\begin{tabular}{lrrrrr}
\toprule
Model & $r$ & RMSE & Rel.\ RMSE & SNR (dB) & Coherence \\
\midrule
\textbf{Closed-form} & \textcolor{good}{\textbf{0.887}} & \textcolor{good}{\textbf{0.474}} & \textcolor{good}{\textbf{0.473}} & \textcolor{good}{\textbf{6.51}} & \textcolor{good}{\textbf{0.715}} \\
Linear SGD & \textcolor{bad}{$-0.652$} & \textcolor{bad}{1.591} & \textcolor{bad}{1.588} & \textcolor{bad}{$-4.02$} & 0.288 \\
FIR Filter & \textcolor{bad}{$-0.210$} & \textcolor{bad}{1.520} & \textcolor{bad}{1.517} & \textcolor{bad}{$-3.62$} & 0.059 \\
Conv Encoder & \textcolor{bad}{$-0.043$} & \textcolor{bad}{1.444} & \textcolor{bad}{1.441} & \textcolor{bad}{$-3.17$} & 0.003 \\
\bottomrule
\end{tabular}
\end{table}

\subsection{Band Power Correlations}

\begin{table}[H]
\centering
\caption{Mean band power correlations per model.}
\label{tab:bandpower}
\begin{tabular}{lrrrrr}
\toprule
Model & $\delta$ & $\theta$ & $\alpha$ & $\beta$ & $\gamma$ \\
\midrule
Closed-form & 0.633 & 0.638 & 0.492 & 0.563 & 0.368 \\
Linear SGD  & 0.258 & 0.240 & 0.129 & 0.277 & 0.131 \\
FIR Filter  & $-0.227$ & $-0.094$ & 0.090 & $-0.207$ & $-0.110$ \\
Conv Encoder & --- & --- & --- & --- & --- \\
\bottomrule
\end{tabular}
\end{table}

\subsection{Loss Ablation Study}

The ablation was run with the Conv Encoder architecture (100 epochs each):

\begin{table}[H]
\centering
\caption{Loss ablation on Conv Encoder.}
\label{tab:ablation_loss}
\begin{tabular}{lrrr}
\toprule
Loss Configuration & $r$ & SNR (dB) & RMSE \\
\midrule
MSE only            & \textbf{0.813} & \textbf{4.71}  & \textbf{0.583} \\
MSE + Spectral      & 0.680          & 2.64           & 0.739 \\
MSE + Band power    & 0.433          & $-0.62$        & 1.076 \\
All losses combined & $-0.341$       & $-4.32$        & 1.649 \\
\bottomrule
\end{tabular}
\end{table}

\textbf{Critical finding:} Each additional loss term \emph{degrades} performance.
The spectral and band-power losses appear to introduce optimization conflicts
or gradient interference.

\subsection{Window Length Ablation}

\begin{table}[H]
\centering
\caption{Window length ablation on Conv Encoder (combined loss, 100 epochs).}
\label{tab:ablation_window}
\begin{tabular}{lrrr}
\toprule
Window $T$ (samples) & $r$ & SNR (dB) & RMSE \\
\midrule
128  & $-0.019$ & $-3.02$ & 1.418 \\
256  & $-0.367$ & $-4.38$ & 1.660 \\
512  & $-0.466$ & $-4.84$ & 1.749 \\
1024 & $-0.011$ & $-3.32$ & 1.468 \\
\bottomrule
\end{tabular}
\end{table}

All window sizes yield negative correlations when using the combined loss,
confirming that the primary issue is the loss function/training dynamics, not
the window size.

\section{Analysis and Diagnosis}

\subsection{Why Gradient-Trained Models Fail}

Several observations point to root causes:

\begin{enumerate}
  \item \textbf{Loss magnitude explosion}: Training losses start in the millions
    (e.g., $3.8 \times 10^6$ for Conv Encoder epoch 1), suggesting the data
    or loss is not properly scaled. The closed-form solution bypasses this
    entirely.

  \item \textbf{Combined loss interference}: The ablation clearly shows that
    adding spectral and band-power losses progressively destroys performance.
    MSE-only achieves $r=0.81$, nearly matching the closed-form baseline.
    The auxiliary losses likely create conflicting gradients.

  \item \textbf{SGD vs.\ analytic}: The linear SGD model has identical architecture
    to the closed-form but achieves $r=-0.65$ vs.\ $r=0.89$. This means the
    Adam optimizer with the given hyperparameters cannot find the analytic
    optimum --- a clear optimization failure, not an architecture limitation.

  \item \textbf{Overfitting to wrong objective}: The Conv Encoder with combined
    loss converges to a solution that is \emph{anti-correlated} with the target,
    suggesting it gets trapped in a bad local minimum where spectral matching
    conflicts with time-domain accuracy.
\end{enumerate}

\subsection{The Closed-Form Solution as Upper Bound}

With $r=0.887$ and coherence of 0.715, the closed-form linear spatial filter
explains approximately $78.7\%$ of variance in the in-ear signal ($r^2$).
This is expected given the linear forward model with 10\,dB SNR. The remaining
$21.3\%$ is attributable to additive sensor noise.

The theoretical SNR limit for a perfect linear decoder with 10\,dB input SNR
on our 4-channel output is approximately $10 + 10\log_{10}(21/4) \approx 17.2$\,dB
(assuming perfect spatial averaging). The achieved 6.51\,dB falls short because
the mixing matrix concentrates on a subset of channels, limiting the effective
spatial averaging gain.

\section{Proposed Improvements for Iteration 002}

Based on our diagnosis, we propose the following changes:

\begin{enumerate}
  \item \textbf{Use MSE-only loss}: The ablation definitively shows this is
    the best configuration. Remove spectral and band-power losses from the
    default configs.

  \item \textbf{Initialize from closed-form solution}: Use $\mathbf{W}^*$ as
    initialization for gradient-trained models. For the linear SGD model, this
    should immediately close the gap. For deeper models, the spatial layer can
    be initialized this way.

  \item \textbf{Fix loss scaling}: Normalize the loss by window length and
    number of channels to bring magnitudes to $\mathcal{O}(1)$. Current losses
    in the millions make learning rate tuning extremely brittle.

  \item \textbf{Add gradient clipping}: Clip gradients to prevent the explosion
    evident in early training epochs.

  \item \textbf{Learning rate warmup}: Add linear warmup for the first 10\%
    of training to stabilize early optimization.

  \item \textbf{Proper hyperparameter sweep}: Run Optuna sweeps with MSE-only
    loss to find optimal learning rates for each architecture.

  \item \textbf{Increase training data}: 2,090 windows may be insufficient
    for the 24K-parameter Conv Encoder. Increase \texttt{n\_subjects} and
    \texttt{n\_samples}.
\end{enumerate}

\section{Conclusion}

This first iteration establishes that the closed-form linear spatial filter
is a strong baseline ($r=0.887$), but gradient-trained models fail due to
training instabilities. The primary issue is the combined loss function and
poor loss scaling, not the model architectures themselves. The MSE-only ablation
proves that the Conv Encoder \emph{can} achieve $r=0.81$ with proper loss
configuration. Iteration~002 will implement the fixes above and target
$r \geq 0.85$ for all gradient-trained models.

\end{document}
